% Part: first-order-logic
% Chapter: syntax-and-semantics
% Section: first-order-languages

\documentclass[../../../include/open-logic-section]{subfiles}

\begin{document}

\olfileid[hi]{fol}{syn}{fol}

\olsection{प्रथम-क्रम भाषाएँ}


प्रथम-क्रम तर्क की अभिव्यक्तियाँ उस मूल शब्दावली से निर्मित होती हैं जिसमें
\emph{!!{variable}s}, \emph{!!{constant}s}, \emph{!!{predicate}s} और कभी-कभी
\emph{!!{function}s} होते हैं। इनसे, तथा तार्किक संयोजकों, परिमाणकों और
कोष्ठक एवं अल्पविराम जैसे विराम-चिह्नों से, \emph{पद} और
\emph{!!{formula}s} बनते हैं।

\begin{explain}
अनौपचारिक रूप से, !!{predicate}s गुणधर्मों और संबंधों के नाम हैं,
!!{constant}s अलग-अलग वस्तुओं के नाम हैं और !!{function}s प्रतिचित्रणों के
नाम हैं। !!{identity}~$\eq$ को छोड़कर ये \emph{अतार्किक प्रतीक} हैं और
मिलकर एक भाषा बनाते हैं। कोई भी प्रथम-क्रम भाषा~$\Lang L$ अपने अतार्किक
प्रतीकों से निर्धारित होती है। सबसे सामान्य स्थिति में~$\Lang L$ में हर
प्रकार के अनंततः अनेक प्रतीक होते हैं।
\end{explain}

सामान्य स्थिति में प्रथम-क्रम तर्क में हम निम्न प्रतीकों का उपयोग करते हैं:

\begin{enumerate}
\item तार्किक प्रतीक
\begin{enumerate}
\item तार्किक संयोजक:
  \startycommalist
  \iftag{prvNot}{\ycomma $\lnot$ (निषेध)}{}%
  \iftag{prvAnd}{\ycomma $\land$ (संयोजन)}{}%
  \iftag{prvOr}{\ycomma $\lor$ (वियोजन)}{}%
  \iftag{prvIf}{\ycomma $\lif$ (!!{conditional})}{}%
  \iftag{prvIff}{\ycomma $\liff$ (!!{biconditional})}{}%
  \iftag{prvAll}{\ycomma $\lforall$ (सार्वत्रिक परिमाणक)}{}%
  \iftag{prvEx}{\ycomma $\lexists$ (अस्तित्वपरक परिमाणक)}{}।
\tagitem{prvFalse}{!!{falsity} के लिए प्रतिज्ञप्ति नियतांक~$\lfalse$।}{}
\tagitem{prvTrue}{!!{truth} के लिए प्रतिज्ञप्ति नियतांक~$\ltrue$।}{}
\item दो-स्थानी !!{identity}~$\eq$।
\item !!{variable}s का एक !!{denumerable}s समुच्चय: $\Obj v_0$, $\Obj v_1$, $\Obj
  v_2$, \dots
\end{enumerate}
\item प्रथम-क्रम तर्क की \emph{मानक भाषा} बनाने वाले अतार्किक प्रतीक
\begin{enumerate}
\item $n$-स्थानी !!{predicate}s का एक !!{denumerable}s समुच्चय, प्रत्येक $n>0$ के लिए: $\Obj
  A^n_0$, $\Obj A^n_1$, $\Obj A^n_2$, \dots
\item !!{constant}s का एक !!{denumerable}s समुच्चय: $\Obj c_0$, $\Obj c_1$, $\Obj
  c_2$, \dots.
\item $n$-स्थानी !!{function}s का एक !!{denumerable}s समुच्चय, प्रत्येक $n>0$ के लिए:
  $\Obj f^n_0$, $\Obj f^n_1$, $\Obj f^n_2$, \dots
\end{enumerate}
\item विराम-चिह्न: (, ), और अल्पविराम।
\end{enumerate}

हमारी अधिकांश परिभाषाएँ और परिणाम प्रथम-क्रम तर्क की पूरी मानक भाषा के लिए
दिए जाएँगे। किंतु अनुप्रयोग के अनुसार हम भाषा को केवल कुछ !!{predicate}s,
!!{constant}s और !!{function}s तक सीमित भी कर सकते हैं।

\begin{ex}
अंकगणित की भाषा~$\Lang L_A$ में एक दो-स्थानी !!{predicate}~$<$, एक
!!{constant}~$\Obj 0$, एक एक-स्थानी !!{function}~$\prime$, और दो दो-स्थानी
!!{function}s~$+$ तथा~$\times$ होते हैं।
\end{ex}

\begin{ex}
समुच्चय सिद्धांत की भाषा~$\Lang L_Z$ में केवल एक दो-स्थानी
!!{predicate}~$\in$ होता है।
\end{ex}

\begin{ex}
क्रमों की भाषा~$\Lang L_\le$ में केवल दो-स्थानी !!{predicate}~$\le$ होता है।
\end{ex}

ये भी रूढ़ियाँ ही हैं: औपचारिक रूप से ये केवल उपनाम हैं। उदाहरणतः $<$,
$\in$ और $\le$, $\Obj A^2_0$ के उपनाम हैं; $\Obj 0$, $\Obj c_0$ का;
$\prime$, $\Obj f^1_0$ का; $+$, $\Obj f^2_0$ का; और $\times$,
$\Obj f^2_1$ का उपनाम है।

\iftag{defNot,defOr,defAnd,defIf,defIff,defTrue,defFalse,defEx,defAll}{%
ऊपर प्रस्तुत आदिम संयोजकों और
\iftag{notprvEx,notprvAll}{परिमाणक}{परिमाणकों} के अतिरिक्त हम निम्नलिखित
\emph{परिभाषित} प्रतीकों का भी उपयोग करते हैं:
\startycommalist
  \iftag{defNot}{\ycomma $\lnot$ (निषेध)}{}%
  \iftag{defAnd}{\ycomma $\land$ (संयोजन)}{}%
  \iftag{defOr}{\ycomma $\lor$ (वियोजन)}{}%
  \iftag{defIf}{\ycomma $\lif$ (!!{conditional})}{}%
  \iftag{defIff}{\ycomma $\liff$ (!!{biconditional})}{}%
  \iftag{defAll}{\ycomma $\lforall$ (सार्वत्रिक परिमाणक)}{}%
  \iftag{defEx}{\ycomma $\lexists$ (अस्तित्वपरक परिमाणक)}{}%
  \iftag{defFalse}{\ycomma !!{falsity}~$\lfalse$}{}%
  \iftag{defTrue}{\ycomma !!{truth}~$\ltrue$}}{}.

\begin{tagblock}{defNot,defOr,defAnd,defIf,defIff,defTrue,defFalse,defEx,defAll}
\begin{explain}
परिभाषित प्रतीक औपचारिक रूप से भाषा का भाग नहीं होता, बल्कि उसे अनौपचारिक
संक्षेप के रूप में प्रस्तुत किया जाता है: उससे हम उन सूत्रों को संक्षिप्त लिख
सकते हैं जो केवल आदिम प्रतीकों के प्रयोग पर बहुत लंबे हो जाते। यह स्पष्टतः
लाभदायक है। किंतु बड़ा लाभ यह है कि प्रमाण छोटे हो जाते हैं। यदि कोई प्रतीक
आदिम हो, तो प्रमाणों में उसे अलग से संभालना पड़ता है। अतः आदिम प्रतीक जितने
अधिक होंगे, हमारे प्रमाण उतने ही लंबे होंगे।
\end{explain}
\end{tagblock}

% Alternate symbols

\begin{intro}
हो सकता है कि आप ऊपर प्रयुक्त पदावली और प्रतीकों से भिन्न रूपों से परिचित
हों। तर्कशास्त्र की पुस्तकों (और शिक्षकों) में ``निषेध'' के लिए प्रायः
$\sim$, $\neg$, और~!\ तथा ``संयोजन'' के लिए $\wedge$, $\cdot$, और $\&$
प्रयुक्त होते हैं। ``आपादन'' या ``निहितार्थ'' के लिए सामान्यतः
$\rightarrow$, $\Rightarrow$, और $\supset$ प्रतीक प्रयुक्त होते हैं।
\iftag{prvIff,defIff}{``द्विशर्त,'' ``द्वि-आपादन,'' या ``(शाब्दिक)
  तुल्यता'' के प्रतीक $\leftrightarrow$,
  $\Leftrightarrow$, और $\equiv$ हैं।}{}
\iftag{prvFalse,defFalse}{$\lfalse$ प्रतीक को अलग-अलग प्रसंगों में
  ``असत्यता,'' ``असत्य,'' ``असंगति,'' या ``तल'' कहा जाता है।}{}
\iftag{prvTrue,defTrue}{$\ltrue$ प्रतीक को अलग-अलग प्रसंगों में
  ``सत्यता,'' ``सत्य,'' या ``शीर्ष'' कहा जाता है।}{}

रूढ़ि के अनुसार !!{constant}s (जिन्हें कभी-कभी नाम कहा जाता है) के लिए
लैटिन वर्णमाला के आरंभ के छोटे अक्षर (जैसे $a$, $b$, $c$) और !!{variable}s
के लिए अंत के छोटे अक्षर (जैसे $x$, $y$, $z$) प्रयुक्त होते हैं। परिमाणक
!!{variable}s, जैसे $x$, के साथ जुड़ते हैं। सार्वत्रिक परिमाणक के संकेतन के
रूपांतरों में $\forall x$, $(\forall x)$, $(x)$, $\Pi x$, $\bigwedge_x$ और
अस्तित्वपरक परिमाणक के रूपांतरों में $\exists x$, $(\exists
x)$, $(Ex)$,
$\Sigma x$, $\bigvee_x$ सम्मिलित हैं।
\end{intro}

\begin{explain}
हम सभी प्रतिज्ञप्ति संक्रियकों और दोनों परिमाणकों को भाषा के आदिम प्रतीक मान
सकते हैं। इसके बजाय हम आदिम प्रतीकों का छोटा भंडार चुनकर शेष !!{operator}s
को परिभाषित भी मान सकते हैं। बूल संक्रियकों के ``सत्य-फलनीय रूप से पूर्ण''
समुच्चयों में $\{ \lnot, \lor \}$, $\{ \lnot, \land \}$, और $\{ \lnot,
\lif\}$ सम्मिलित हैं---अभिव्यंजक रूप से पूर्ण प्रथम-क्रम भाषा पाने के लिए
इनमें से किसी को किसी भी एक परिमाणक के साथ जोड़ा जा सकता है।

आप दो अन्य !!{operator}s से भी परिचित हो सकते हैं: हेनरी शेफ़र के नाम पर
रखा गया शेफ़र स्ट्रोक~$|$, और पियर्स तीर~$\downarrow$, जिसे क्वाइन का डैगर
भी कहते हैं। क्रमशः ``नैन्ड'' और ``नॉर'' के अपने सामान्य पाठ के साथ ये
संक्रियक अपने-आप में सत्य-फलनीय रूप से पूर्ण हैं।
\end{explain}

\end{document}
